\section{Arrays and array lists}

\solution{Problem}{1}
\idea Read the string run by run: remember the current character, count how far it
repeats, then write the character and the count. A count can have several digits; find
how many by repeated division by 10, then fill the digits \emph{from the right}, because
\cc{run \% 10} produces the last digit first. \cc{compress} works in place because the
write position never overtakes the read position: a run of length $r\ge2$ occupies $r$
cells and is replaced by $1+\lfloor\log_{10}r\rfloor+1\le r$ cells.
\listingfile{sol/p01.c}{Run-length encoding, in-place compression and decoding.}
\complexity $\Theta(n)$ time: each input character is read once and each output character
written once. \cc{compress} uses $O(1)$ extra space.
\mistakes Writing the count as a single character (\cc{'0'+run}), which breaks for runs of
10 or more; writing the digits in reverse; forgetting the terminating \cc{'\textbackslash
0'} in \cc{rleEncode}; in \cc{compress}, writing the count \cc{1} for single
characters, or testing \cc{a[r]==c} without \cc{r<n}.

\solution{Problem}{2}
\idea $t_{ji}=a_{ij}$ is a direct double loop. In place, swap $a_{ij}$ and $a_{ji}$ for
the cells \emph{above the diagonal} only ($j>i$): visiting every pair twice swaps each
pair back and leaves the matrix unchanged.
\listingfile{sol/p02.c}{Transpose.}
\complexity $\Theta(nm)$ and $\Theta(n^2)$ time; $O(1)$ extra space in place.
\mistakes Looping $j$ from 0 in the in-place version (double swap); mixing up $n$ and $m$
in the bounds of the rectangular case.

\solution{Problem}{3}
\idea The sum $1+2+\dots+(m+1)=(m+1)(m+2)/2$ is known; subtracting the sum of the array
leaves the missing number. For the bonus, note that before the gap $a[i]=i+1$ and from the
gap on $a[i]=i+2$: the predicate ``$a[i]\ne i+1$'' is false\ldots false true\ldots true, so
bisection finds the first index where it holds.
\listingfile{sol/p03.c}{The missing number, linear and logarithmic.}
\complexity $\Theta(m)$ and $\Theta(\log m)$ time; $O(1)$ space.
\mistakes Using $m(m+1)/2$ (the range has $m+1$ numbers); in the bisection, initialising
\cc{hi=m-1}, which misses the case where the missing number is $m+1$.

\solution{Problem}{4}
\idea \cc{lowerBound} and \cc{upperBound} are the bisection of Section~\ref{sec:bsearch}
with \cc{<} and \cc{<=}. Then \cc{search} is a lower bound plus one check, \cc{count} is
the difference of the two bounds, \cc{insert} opens a gap at the upper bound (so equal
items keep their arrival order), and \cc{erase} closes the gap at the lower bound.
\listingfile{sol/p04.c}{A sorted array list.}
\complexity Bounds, \cc{search}, \cc{count}: $\Theta(\log n)$. \cc{insert}, \cc{erase}:
$\Theta(\log n)$ comparisons plus up to $n$ moves, $\Theta(n)$.
\mistakes Reading \cc{items[p]} when $p=\cc{size}$ (not found and larger than every
item); shifting in the wrong direction (left-to-right when opening a gap overwrites every
item with the same value); counting by a linear scan from the lower bound, which is
$\Theta(n)$ when all items are equal.

\solution{Problem}{5}
\idea Reduce $k$ modulo $n$. Rotating right by $k$ moves the last $k$ items to the front.
Reverse the whole array, then reverse the first $k$ items and the remaining $n-k$ items
separately; each block is now in its original order but the blocks have swapped places.
\listingfile{sol/p05.c}{Rotation by three reversals.}
\complexity $\Theta(n)$ time (each item is swapped at most twice), $O(1)$ space.
\mistakes Forgetting \cc{k \% n} ($k$ may be $10^9$); computing \cc{k \% n} when $n=0$
(division by zero); rotating left instead of right (reversing the blocks in the other
order).

\solution{Problem}{6}
\idea Two indices walk through \cc{a} and \cc{b}; at each step the smaller front item is
copied. Comparing with \cc{<=} copies \cc{a}'s item on ties (stability). When one array
is exhausted, the rest of the other is copied.
\listingfile{sol/p06.c}{Merging two sorted arrays.}
\complexity $\Theta(n+m)$: each item is copied once.
\mistakes Forgetting the two tail loops; using \cc{<} (the order of equal items is then
\cc{b} first, which the statement forbids).

\solution{Problem}{7}
\idea Keep a write index $w$: \cc{a[0..w-1]} holds the distinct values seen so far. Since
the array is sorted, a value is new exactly when it differs from \cc{a[w-1]}.
\listingfile{sol/p07.c}{Remove duplicates from a sorted array.}
\complexity $\Theta(n)$ time, $O(1)$ space.
\mistakes Comparing with \cc{a[r-1]} instead of \cc{a[w-1]} (works here, but not for the
``at most two copies'' variant); not handling $n=0$.

\solution{Problem}{8}
\idea Two pointers from both ends. If the sum is too small, only a larger left item can
help: \cc{lo++}. If it is too large, \cc{hi--}. No pair is skipped: when \cc{lo} moves,
\cc{a[lo]} could only pair with indices $\le\cc{hi}$, all of which give too small a sum.
\listingfile{sol/p08.c}{Two-sum in a sorted array.}
\complexity $\Theta(n)$ time, $O(1)$ space (the double loop is $\Theta(n^2)$).
\mistakes The condition \cc{lo<=hi}, which may use the same index twice; moving both
pointers at once.

\solution{Problem}{9}
\idea Kadane's algorithm. Let $E_i$ be the best sum of a subarray \emph{ending} at $i$.
Either it extends the best one ending at $i-1$, or it starts afresh at $i$:
$E_i=a_i+\max(E_{i-1},0)$. The answer is $\max_iE_i$.
\listingfile{sol/p09.c}{Maximum subarray sum, $\Theta(n)$ and $\Theta(n^2)$.}
\complexity $\Theta(n)$ time, $O(1)$ space.
\mistakes Initialising \cc{best = 0}, which returns 0 for an all-negative array (the
subarray must be non-empty); resetting to 0 instead of to \cc{a[i]}.

\solution{Problem}{10}
\idea Copy every non-zero item to the next write position (a stable compaction), then fill
the rest with zeros.
\listingfile{sol/p10.c}{Move zeroes.}
\complexity $\Theta(n)$ time, $O(1)$ space.
\mistakes Swapping a zero with the \emph{last} item, which destroys the order; using a
second array.

\solution{Problem}{11}
\idea Boyer--Moore voting. Pair off two different items and discard both: the majority
item, which occurs more than half the time, cannot be paired off completely, so it
survives. The variable \cc{votes} is the number of unpaired copies of \cc{candidate}.
\listingfile{sol/p11.c}{Majority element.}
\complexity $\Theta(n)$ time, $O(1)$ space; the counting version is $\Theta(n^2)$.
\mistakes Assuming the most frequent item is the answer when no majority is guaranteed
(then a second pass must verify the candidate).

\solution{Problem}{12}
\idea A clockwise rotation sends $(i,j)$ to $(j,n-1-i)$. That is the transpose
$(i,j)\mapsto(j,i)$ followed by the reflection of each row $(j,i)\mapsto(j,n-1-i)$; both
are in-place swaps.
\listingfile{sol/p12.c}{Rotating a matrix by $90^\circ$ clockwise.}
\complexity $\Theta(n^2)$ time, $O(1)$ space.
\mistakes Reversing the columns instead of the rows (counter-clockwise rotation);
transposing with $j$ from 0.
